\documentclass[11pt]{article}

\usepackage[a4paper,margin=1in]{geometry}
\usepackage{amsmath,amssymb,amsthm,booktabs}
\usepackage{microtype}
\usepackage{url}

\newtheorem{theorem}{Theorem}
\newtheorem{lemma}[theorem]{Lemma}

\title{Fisher-Strengthened Self-Amalgam Bounds and Partial Ranges for\\
Atlas 153, 171, 174, and 188}
\author{}
\date{}

\begin{document}

\maketitle

\section{Result}

For a finite graph $H$ and a graphon $W$, write
\[
 p=t(K_2,W),
 \qquad
 \Phi_H(p)=(1-p)^{v(H)}
 \chi_H\!\left(\frac1{1-p}\right),
\]
with the usual polynomial continuation at $p=1$.

Put
\[
 \pi(y)=\frac{(1-y)(1+3y)}2,
 \qquad
 g(y)=\frac32y(1-y)^2.
\]
For each of the four graphs, let $\eta_H$ be the unique zero in
$[0,1/3]$ of the integer polynomial $C_H$ displayed in
Section~\ref{sec:fisher-comparison}, and put $\beta_H=\pi(\eta_H)$.

\begin{theorem}[Improved arbitrary-graphon ranges]\label{thm:main}
For every graphon $W$, Atlas $H$ satisfies
\[
 t(H,W)\geq\Phi_H(p)
 \qquad\left(\frac12\leq p\leq\beta_H\right)
\]
for $H=153,171,174,188$.  The exact endpoints, specified as roots of
the displayed integer polynomials, have the certified brackets
\[
\begin{array}{c|c|c}
H&\text{certified bracket for }\beta_H&
100\,\dfrac{\beta_H-1/2}{1-1/2}\\ \hline
153&(0.6181289270058,0.6181289270059)&23.6257854\ldots\%\\
171&(0.6309517719656,0.6309517719657)&26.1903543\ldots\%\\
174&(0.5934675772766,0.5934675772767)&18.6935154\ldots\%\\
188&(0.6139911259451,0.6139911259452)&22.7982251\ldots\%.
\end{array}
\]
Thus the safe rounded-down closed endpoints are respectively
$0.6181289270$, $0.6309517719$, $0.5934675772$, and $0.6139911259$.
\end{theorem}

The proof is not a finite-step reduction.  The elementary fallback uses
weighted Cauchy--Schwarz and the already proved pure-chordal bounds for the
diamond and Atlas $47$.  The improvement uses, in addition, Fisher's sharp
$k=2$ triangle-density band.  This is the accepted, Lean-formalized
triangle theorem already used for Atlas $148$; it is not a flag-algebra
input.

\section{A self-amalgam inequality over an arbitrary root}

Let $J$ be the subgraph of a finite graph $F$ induced by a chosen set of
root vertices.  Form $S=F\cup_JF$ from two labelled copies of $F$ by
identifying their copies of $J$ and adding no edges between the two
branches.  The graph $J$ need not be a clique.

\begin{lemma}[Two-copy self-amalgamation]\label{lem:amalgam}
For every graphon $W$,
\[
 t(S,W)t(J,W)\geq t(F,W)^2.
\]
If a graph $H$ is a spanning subgraph of $S$, then also
$t(H,W)\geq t(S,W)$.
\end{lemma}

\begin{proof}
Fix the variables $x_J$ on the common root.  Let
\[
 h(x_J)=\prod_{ab\in E(J)}W(x_a,x_b)
\]
and let $g(x_J)$ be the integral of all edge factors of $F$ not already
in $J$, over the variables in $V(F)\setminus V(J)$.  Tonelli--Fubini
gives
\[
 t(J,W)=\int h,
 \qquad
 t(F,W)=\int hg,
 \qquad
 t(S,W)=\int hg^2.
\]
Weighted Cauchy--Schwarz, applied to $\sqrt h$ and $\sqrt h\,g$, gives
\[
 \left(\int hg\right)^2
 \leq\left(\int h\right)\left(\int hg^2\right).
\]
This form uses no division and therefore also covers $t(J,W)=0$.

Finally, $0\leq W\leq1$.  Removing the edge factors in
$E(S)\setminus E(H)$ can only increase the integrand, so a spanning
subgraph $H\subseteq S$ satisfies $t(H,W)\geq t(S,W)$.
\end{proof}

We shall also need an upper bound, rather than Sidorenko's lower bound,
for the four-vertex path $P_4=abcd$:
\begin{equation}\label{eq:p4-upper}
\begin{split}
 t(P_4,W)
 &=\int W(x_a,x_b)W(x_b,x_c)W(x_c,x_d)\,d\mu^4\\
 &\leq\int W(x_a,x_b)W(x_c,x_d)\,d\mu^4=p^2.
\end{split}
\end{equation}
This is pointwise deletion of the middle edge.  In particular, no
division by $t(P_4,W)$ will occur below.

\section{The two positive bases}

Let $D$ be the diamond, NetworkX Atlas $17$, and let $B$ be NetworkX
Atlas $47$, and put $T=t(K_3,W)$.  Their edge lists and chromatic
polynomials are
\[
\begin{array}{c|l|l}
F&E(F)&\chi_F(x)\\ \hline
D&01,02,03,12,23&x(x-1)(x-2)^2\\
B&01,04,12,13,14,23,34&x(x-1)(x-2)^3.
\end{array}
\]
Both graphs are $3$-pure chordal.  The pure-chordal theorem therefore
gives, for every graphon and every $p\geq1/2$,
\begin{equation}\label{eq:base-bounds}
 t(D,W)\geq p(2p-1)^2,
 \qquad
 t(B,W)\geq p(2p-1)^3.
\end{equation}
These right sides are exactly $\Phi_D(p)$ and $\Phi_B(p)$.

We shall also use the triangle-sensitive form of the accepted
pure-chordal entropy-gluing inequality.  The diamond has two maximal
triangles and one edge separator, while $B$ has three maximal triangles
in a path and two edge separators.  Consequently
\begin{equation}\label{eq:triangle-sensitive-bases}
 t(D,W)p\geq T^2,
 \qquad
 t(B,W)p^2\geq T^3.
\end{equation}
Equivalently, because our range has $p\geq1/2$,
\[
 t(D,W)\geq\frac{T^2}{p},
 \qquad
 t(B,W)\geq\frac{T^3}{p^2}.
\]
The first inequality is also Lemma~\ref{lem:amalgam} applied to two
triangles glued along an edge.  The second is the three-clique-tree case
of the same accepted pure-chordal entropy theorem.

\section{Exact Atlas identifications}

All labels in this section are the NetworkX Atlas labels.  We exhibit a
supergraph $S_H\supseteq H$ that is a two-copy self-amalgam of one of the
two bases.

\paragraph{Atlas 153.}
Its edge set is
\[
 01,04,05,12,15,23,25,34.
\]
Add the two edges $03$ and $24$.  With common root $J=\{0,2\}$ and
branches $\{1,5\}$ and $\{3,4\}$, the resulting graph is two diamonds
glued along the nonedge $02$.  Explicit maps from the two rooted copies
to the displayed Atlas-$17$ labelling are
\[
 (1,0,5,2)\mapsto(0,1,2,3),
 \qquad
 (3,0,4,2)\mapsto(0,1,2,3).
\]
They agree on the common root.  Since $J$ has no edges, $t(J,W)=1$.
Lemma~\ref{lem:amalgam} and \eqref{eq:base-bounds} give
\begin{equation}\label{eq:153-lower}
 t(153,W)\geq t(D,W)^2
 \geq p^2(2p-1)^4.
\end{equation}
Using \eqref{eq:triangle-sensitive-bases} instead gives the stronger
triangle-sensitive estimate
\begin{equation}\label{eq:153-triangle-lower}
 t(153,W)\geq\frac{T^4}{p^2}.
\end{equation}

\paragraph{Atlas 171.}
Its edge set is
\[
 01,04,05,12,23,25,34,35,45.
\]
Add $02$ and $13$.  The common root
$J=\{1,2,4,5\}$ has edges $12,25,45$, hence is a four-vertex path;
the branches are the single vertices $0$ and $3$.  The maps
\[
 (1,0,4,5,2)\mapsto(0,1,2,3,4),
 \qquad
 (1,3,4,5,2)\mapsto(0,1,2,3,4)
\]
identify both rooted branches with Atlas $47$ and agree on $J$.

\paragraph{Atlas 174.}
Its edge set is
\[
 01,02,03,14,15,23,25,34,45.
\]
Add $05$ and $24$.  Now $J=\{1,2,3,5\}$ has path edges $15,25,23$,
and the branches are $0$ and $4$.  The two rooted Atlas-$47$ maps are
\[
 (1,0,3,2,5)\mapsto(0,1,2,3,4),
 \qquad
 (1,4,3,2,5)\mapsto(0,1,2,3,4).
\]

\paragraph{Atlas 188.}
Its edge set is
\[
 01,02,05,12,13,15,23,24,34,45.
\]
Add $04$.  Here $J=\{0,2,3,5\}$ has path edges $05,02,23$, and the
branches are $1$ and $4$.  The rooted Atlas-$47$ maps are
\[
 (3,1,5,0,2)\mapsto(0,1,2,3,4),
 \qquad
 (3,4,5,0,2)\mapsto(0,1,2,3,4).
\]

For each of the final three graphs, Lemma~\ref{lem:amalgam},
\eqref{eq:p4-upper}, and \eqref{eq:base-bounds} yield directly
\[
 t(B,W)^2
 \leq t(S_H,W)t(P_4,W)
 \leq t(S_H,W)p^2.
\]
Since $p\geq1/2$, division is only by the positive scalar $p^2$; no
rooted density is divided out.  Therefore
\begin{equation}\label{eq:three-lower}
 t(H,W)\geq t(S_H,W)
 \geq\frac{t(B,W)^2}{p^2}
 \geq(2p-1)^6
 \quad(H=171,174,188).
\end{equation}
The triangle-sensitive base bound in
\eqref{eq:triangle-sensitive-bases} yields at the same time
\begin{equation}\label{eq:three-triangle-lower}
 t(H,W)\geq\frac{T^6}{p^6}
 \quad(H=171,174,188).
\end{equation}

\section{Elementary fallback comparison}\label{sec:elementary-comparison}

Define
\[
\begin{array}{c|l}
H&R_H(p)\\ \hline
153&8p^3-11p^2+7p-2\\
171&32p^4-59p^3+41p^2-12p+1\\
174&32p^4-61p^3+44p^2-13p+1\\
188&32p^4-65p^3+48p^2-14p+1.
\end{array}
\]

The four targets from the target-sign note are
\[
 \Phi_H(p)=p(2p-1)Q_H(p),
\]
where
\[
\begin{array}{c|l}
H&Q_H(p)\\ \hline
153&7p^3-12p^2+8p-2\\
171&11p^3-20p^2+13p-3\\
174&13p^3-25p^2+17p-4\\
188&17p^3-33p^2+22p-5.
\end{array}
\]
Exact expansion of \eqref{eq:153-lower} and
\eqref{eq:three-lower} gives
\begin{align}
 p^2(2p-1)^4-\Phi_{153}(p)
 &=p(p-1)(2p-1)R_{153}(p),\label{eq:gap153}\\
 (2p-1)^6-\Phi_H(p)
 &=(p-1)(2p-1)R_H(p)
 \quad(H=171,174,188).\label{eq:gapthree}
\end{align}

Every $R_H$ is strictly increasing on $[1/2,1]$.  For $H=153$,
$R'_{153}$ is a positive-leading quadratic of discriminant $-188$.
For $H=171,174,188$, the cubic $R'_H$ has respectively discriminant
\[
 -233324,\qquad-669356,\qquad-651240.
\]
Thus it has one real zero; because its leading coefficient is positive
and
\[
 R'_{171}(1/2)=\frac34,
 \qquad
 R'_{174}(1/2)=R'_{188}(1/2)=\frac54,
\]
that zero lies below $1/2$, and $R'_H>0$ thereafter.  Finally,
\[
 R_{153}(1/2)=-\frac14,
 \qquad
 R_{171}(1/2)=R_{174}(1/2)=R_{188}(1/2)=-\frac18,
\]
whereas the four values at $1$ are $2,3,3,2$.  This proves existence
and uniqueness of a zero $\gamma_H$ of each $R_H$ and shows
$R_H\leq0$ on $[1/2,\gamma_H]$.  Equations
\eqref{eq:gap153}--\eqref{eq:gapthree} therefore give the independent,
elementary fallback endpoints
\[
 0.6128767076,\quad 0.6292381397,\quad
 0.5927496223,\quad 0.6121757702
\]
for Atlas $153,171,174,188$, respectively.

For completeness, the decimal brackets are certified by substituting
their finite-decimal rational endpoints.  After reducing to a fraction
with positive denominator, the exact numerator pairs (left, right) are
\[
\begin{array}{c|r|r}
H&\operatorname{num}R_H(\text{left})&
\operatorname{num}R_H(\text{right})\\ \hline
153&-326087393309890391&10774669858388085533\\
171&-23637245921617628887830851119&1268914984346577868385927201\\
174&-30483718352770372806343598559&4074137222454462766319\\
188&-354905520499384947188549599&27334051028435931971959991681.
\end{array}
\]
The signs are therefore exact, not floating-point evidence.

\section{Fisher strengthening and the sharp certified endpoints}
\label{sec:fisher-comparison}

We use the following already accepted triangle-density input.  If
$1/2<p\leq2/3$, there is a unique $y\in(0,1/3]$ such that
\begin{equation}\label{eq:fisher-parameter}
 p=\pi(y)=\frac{(1-y)(1+3y)}2,
\end{equation}
and Fisher's sharp $k=2$ band gives
\begin{equation}\label{eq:fisher-bound}
 T=t(K_3,W)\geq g(y)=\frac32y(1-y)^2.
\end{equation}
At $p=1/2$ the same inequality follows simply from $T\geq0=g(0)$.
The map $\pi$ is strictly increasing on $[0,1/3]$ because
$\pi'(y)=1-3y$.

This is precisely the triangle theorem used and audited in
\path{notes/atlas148_paw_bias_hilbert_projection.tex}.  In the Lean
development it is \path{fisher_triangle_bound} in
\path{Methods/TriangleDensity.lean}.
Its imported analytic declaration is
\[
 \texttt{triangleDensityLowerBound\_twoThirds}.
\]
It is provided by \path{lean/Taeyoung/Fisher/}.  That development formalizes Fisher's
finite/spectral proof and its graphon sampling bridge.  It uses no flag
algebra.  Its axiom audit is
\[
 \texttt{[propext, Classical.choice, Quot.sound]}.
\]

The four primitive integer comparison polynomials are
\begin{align*}
C_{153}(y)={}&-15309y^{10}+25515y^9-7614y^8-5832y^7-702y^6
 +4266y^5\\
 &{}-1524y^4+248y^3-69y^2-5y+2,\\[2mm]
C_{171}(y)={}&-1948617y^{14}+649539y^{13}+1968300y^{12}
 -354294y^{11}-923643y^{10}\\
 &{}+21141y^9+83106y^8+173988y^7-64071y^6+4437y^5\\
 &{}-2640y^4-494y^3+11y^2+19y+2,\\[2mm]
C_{174}(y)={}&-2302911y^{14}+767637y^{13}+2165130y^{12}
 -472392y^{11}-949887y^{10}\\
 &{}+82377y^9+86022y^8+161352y^7-63261y^6+6039y^5\\
 &{}-3282y^4-1024y^3-101y^2+11y+2,\\[2mm]
C_{188}(y)={}&-3011499y^{14}+1003833y^{13}+2794986y^{12}
 -629856y^{11}-1159839y^{10}\\
 &{}+152361y^9+126846y^8+136080y^7-75249y^6+8595y^5\\
 &{}-1002y^4-560y^3-69y^2+11y+2.
\end{align*}

Substitute \eqref{eq:fisher-parameter} into the targets and use the lower
bounds \eqref{eq:153-triangle-lower} and
\eqref{eq:three-triangle-lower}.  Exact expansion gives
\begin{align}
 \pi(y)^2\left(\frac{g(y)^4}{\pi(y)^2}
       -\Phi_{153}(\pi(y))\right)
 &=\frac{y(1-y)^3}{64}C_{153}(y),
 \label{eq:fisher-gap153}\\
 \pi(y)^6\left(\frac{g(y)^6}{\pi(y)^6}
       -\Phi_H(\pi(y))\right)
 &=\frac{y(1-y)^7}{1024}C_H(y),
 \quad H=171,174,188.
 \label{eq:fisher-gapthree}
\end{align}
All factors outside $C_H$ are nonnegative on the parameter interval, and
the powers of $\pi(y)$ on the left are positive.

It remains only to locate the first zero of $C_H$.  Use the signed Sturm
sequence
\[
 C_H,\quad C'_H,\quad-\operatorname{rem}(C_H,C'_H),\quad\ldots
\]
and delete zeros before counting sign variations.  Exact rational Euclidean
division gives
\[
\begin{array}{c|c|c|c|c}
H&\deg C_H&V(0)&V(1/3)&V(0)-V(1/3)\\ \hline
153&10&6&5&1\\
171&14&8&7&1\\
174&14&8&7&1\\
188&14&8&7&1.
\end{array}
\]
Thus every $C_H$ has exactly one zero $\eta_H$ in $(0,1/3)$.  Moreover
$C_H(0)=2$, and exact Horner evaluation at the rational finite-decimal
endpoints below has signs $+$ and $-$:
\[
\begin{array}{c|c|c}
H&\text{exact rational bracket for }\eta_H&
  \text{induced exact bracket for }\beta_H=\pi(\eta_H)\\ \hline
153&(0.1534486672074800,0.1534486672074801)&
 (0.6181289270058,0.6181289270059)\\
171&(0.1790286678151776,0.1790286678151777)&
 (0.6309517719656,0.6309517719657)\\
174&(0.1124274870396899,0.1124274870396900)&
 (0.5934675772766,0.5934675772767)\\
188&(0.1459379567956481,0.1459379567956482)&
 (0.6139911259451,0.6139911259452).
\end{array}
\]
The second column is therefore an exact isolating certificate, not a
floating-point root computation.  The last column follows by exact rational
evaluation of the increasing quadratic $\pi$.  Since $C_H$ is positive
from $0$ through its unique zero, equations
\eqref{eq:fisher-gap153}--\eqref{eq:fisher-gapthree} prove
Theorem~\ref{thm:main}.

\section{Reproducibility and scope}

Run
\begin{verbatim}
python experiments/four_self_amalgam_partial_verify.py
\end{verbatim}
from the repository root.  The script uses only small NetworkX Atlas
graphs, exact SymPy polynomials, and Python rational arithmetic.  It checks
the four edge inclusions and rooted isomorphisms, the target and gap
factorizations, the elementary derivative discriminants, the exact Sturm
variation counts, and every rational endpoint sign and parameter conversion.

These bounds improve four low-density partial intervals but do not close a
row.  Above $\beta_H$, the Fisher-enhanced self-amalgam lower bound used here
is smaller than $\Phi_H$; this is a limitation of the certificate, not
evidence that the desired graphon inequality fails there.  Section
\ref{sec:elementary-comparison} remains a completely independent fallback
that avoids the sharp triangle theorem, at the cost of slightly shorter
intervals.

\end{document}
